\cotefile{130}
\scanpage{20}
% archive sheet number: 10 (inverted)
\mspage{-VII- 9 -}\hfill n°383

4. \uline{Schémas de Dynkin}.

Soient $S$ un système de racines réduit, $G$ son groupe de Weyl, $(u_1,\ldots,u_r)$ ses racines simples relatives à une chambre. Posons $n(i,j)=2(u_i\mid u_j)/(u_i\mid u_i)$. D'après le \S~4, n°2, les seuls cas possibles, pour $i\ne j$, sont les suivants (en tenant compte du fait que l'angle de deux racines simples est obtus ; cf.\ \S~3, n°1, prop.2) :

\[
\begin{aligned}
1)\ & n(i,j)=n(\unsure{j|J},i)=0\ ;\ (u_i,u_j)=\handwritten{\pi}/2\ ;\ s_i s_j=-1\\
2)\ & n(i,j)=n(j,i)=-1\ ;\ (u_i,u_j)=2\handwritten{\pi}/3\ ;\ s_i s_j\ \text{d'ordre }3\ ;\ \|u_i\|=\|u_j\|\\
3)\ & n(i,j)=-1,\ n(j,i)=-2\ ;\ (u_i,u_j)=3\handwritten{\pi}/4\ ;\ s_i s_j\ \text{d'ordre }4\ ;\ \|u_i\|=\handwritten{\sqrt{2}}\,\|u_j\|
\end{aligned}
\]
(ou ce qu'on obtient en échangeant $i$ et $j$)
\[
4)\quad n(i,j)=-1,\ n(j,i)=-3\ ;\ (u_i,u_j)=5\handwritten{\pi}/6\ ;\ s_i s_j\ \text{d'ordre }6\ ;\ \|u_i\|=\handwritten{\sqrt{3}}\,\|u_j\|
\]
(ou ce qu'on obtient en échangeant $i$ et $j$).

Nous appellerons \uline{schéma de Dynkin} de $S$ le graphe à $r$ noeuds $1,2,\ldots,r$ dont les simplexes sont les parties $\handwritten{\{}i,j\handwritten{\}}$ à deux éléments telles que $(u_i\mid u_j)\ne0$, chaque noeud $i$ étant affecté du coefficient $\|u_i\|^2$. Le graphe sous-jacent à ce schéma n'est autre que le graphe de $G$. Les remarques précédentes permettent de reconstituer le schéma de Coxeter de $G$ à partir du schéma de Dynkin de $S$ (mais pas inversement).

PROPOSITION 2.- \uline{Soient} $V,V'$ \uline{des espaces hilbertiens réels de dimension finie}, $S$ (resp.\ $S'$) \uline{un système de racines réduit dans} $V$ (resp.\ $V'$), $u_1,\ldots,u_r$ (resp.\ $u'_1,\ldots,u'_{r'}$) \uline{les racines simples de} $S$ (resp.$S'$) \uline{relatives à une chambre}. \uline{S'il existe un isomorphisme} $i\handwritten{\rightsquigarrow}j(i)$ \uline{du schéma de Dynkin de} $S$ \uline{sur celui de} $S'$, \uline{il existe un isomorphisme de} $V$ \uline{sur} $V'$ \uline{transformant} $u_i$ \uline{en} $u'_{j(i)}$ \uline{pour tout} $i$, \uline{et} $S$ \uline{en} $S'$.

En effet les hypothèses et les remarques ci-dessus entraînent que $r=r'$ et que $(u_i\mid u_k)=(u'_{j(i)}\mid u'_{j(k)})$ quels que soient $i$ et $k$. Donc l'application $x_1u_1+\ldots+x_ru_r\handwritten{\rightsquigarrow}x_1u'_{j(1)}+\ldots+x_ru'_{j(r)}$ est un isomorphisme $f$ de $V$ sur $V'$, qui transforme le groupe de Weyl $G$ de $S$ en le groupe de Weyl $G'$ de $S'$. Or $S$ (resp.\ $S'$)
